\documentclass[11pt,a4paper,oneside]{report}

% =================================================
% 1. PACKAGES
% =================================================

% Langue et encodage
\usepackage[T1]{fontenc}
\usepackage[utf8]{inputenc}
\usepackage[french]{babel}
\usepackage{csquotes}
\usepackage{lmodern}
\usepackage{microtype}

% Mise en page
\usepackage[a4paper,margin=2.5cm]{geometry}
\usepackage{setspace}
\usepackage{parskip}

% Mathématiques
\usepackage{amsmath,amssymb,amsthm,mathtools}
\usepackage{siunitx}

% Images et couleurs
\usepackage{graphicx}
\usepackage{xcolor}

% Tableaux
\usepackage{booktabs}
\usepackage{array}
\usepackage{tabularx}

% Code source
\usepackage{listings}

% Encadrés et listes
\usepackage[most]{tcolorbox}
\usepackage{enumitem}

% En-têtes et pieds de page
\usepackage{fancyhdr}

% Bibliographie
\usepackage[
    backend=biber,
    style=numeric
]{biblatex}
\addbibresource{bibliographie/references.bib}

% Hyperliens et références
\usepackage{hyperref}
\usepackage[nameinlink,noabbrev]{cleveref}

% =================================================
% 2. CONFIGURATION GENERALE
% =================================================

\setstretch{1.15}

\hypersetup{
    colorlinks=true,
    linkcolor=blue!60!black,
    citecolor=green!40!black,
    urlcolor=blue!70!black,
    pdftitle={MiniNumpy - Carnet de mathématiques},
    pdfauthor={Louis Subtil}
}

% =================================================
% 3. RACCOURCIS MATHEMATIQUES
% =================================================

% Ensembles
\newcommand{\N}{\mathbb{N}}
\newcommand{\Z}{\mathbb{Z}}
\newcommand{\Q}{\mathbb{Q}}
\newcommand{\R}{\mathbb{R}}
\newcommand{\C}{\mathbb{C}}

% Opérateurs
\DeclareMathOperator{\rg}{rg}
\DeclareMathOperator{\Vect}{Vect}
\DeclareMathOperator{\Card}{Card}

% Notations
\newcommand{\abs}[1]{\left|#1\right|}
\newcommand{\norm}[1]{\left\|#1\right\|}
\newcommand{\vect}[1]{\mathbf{#1}} 

% =================================================
% 4. ENVIRONNEMENTS MATHEMATIQUES
% =================================================

\theoremstyle{definition}

\newtheorem{definition}{Définition}[chapter]
\newtheorem{example}{Exemple}[chapter]

\theoremstyle{plain}

\newtheorem{theorem}[definition]{Théorème}
\newtheorem{proposition}[definition]{Proposition}
\newtheorem{lemma}[definition]{Lemme}
\newtheorem{corollary}[definition]{Corollaire}

\theoremstyle{remark}

\newtheorem{remark}[definition]{Remarque}

% =================================================
% 5. ENCADRES PERSONNALISES
% =================================================

\newtcolorbox{importantbox}{
    breakable,
    colback=red!5,
    colframe=red!60!black,
    title=À retenir,
    fonttitle=\bfseries
}

\newtcolorbox{methodbox}{
    breakable,
    colback=green!5,
    colframe=green!50!black,
    title=Méthode,
    fonttitle=\bfseries
}

\newtcolorbox{codebox}{
    breakable,
    colback=gray!5,
    colframe=gray!60,
    title=Implémentation informatique,
    fonttitle=\bfseries
}

% =================================================
% 6. CONFIGURATION DU CODE SOURCE
% =================================================

\definecolor{codegreen}{RGB}{0,128,0}
\definecolor{codegray}{RGB}{100,100,100}
\definecolor{codepurple}{RGB}{150,50,150}
\definecolor{codebackground}{RGB}{245,245,245}

\lstdefinestyle{mystyle}{
    backgroundcolor=\color{codebackground},
    commentstyle=\color{codegreen},
    keywordstyle=\color{blue},
    stringstyle=\color{codepurple},
    numberstyle=\tiny\color{codegray},
    basicstyle=\ttfamily\small,
    numbers=left,
    numbersep=8pt,
    breaklines=true,
    showstringspaces=false,
    frame=single,
    tabsize=4,
    captionpos=b,
    columns=fullflexible,
    literate=
        {é}{{\'e}}1 {è}{{\`e}}1 {ê}{{\^e}}1 {ë}{{\"e}}1
        {à}{{\`a}}1 {â}{{\^a}}1 {ä}{{\"a}}1
        {î}{{\^i}}1 {ï}{{\"i}}1
        {ô}{{\^o}}1 {ö}{{\"o}}1
        {ù}{{\`u}}1 {û}{{\^u}}1 {ü}{{\"u}}1
        {ç}{{\c{c}}}1 {œ}{{\oe}}1 {æ}{{\ae}}1
        {É}{{\'E}}1 {È}{{\`E}}1 {Ê}{{\^E}}1 {À}{{\`A}}1
        {Ç}{{\c{C}}}1 {Œ}{{\OE}}1
        {«}{{\guillemotleft}}1 {»}{{\guillemotright}}1
        {°}{{\textdegree}}1 {€}{{\texteuro}}1
}

\lstset{style=mystyle}

% =================================================
% 7. EN-TETES ET PIEDS DE PAGE
% =================================================

\pagestyle{fancy}
\fancyhf{}

\fancyhead[L]{\nouppercase{\leftmark}}
\fancyhead[R]{MiniNumpy}

\fancyfoot[C]{\thepage}

\setlength{\headheight}{15pt}

% Appliquer également le style aux pages de chapitres
\fancypagestyle{plain}{
    \fancyhf{}
    \fancyhead[L]{\nouppercase{\leftmark}}
    \fancyhead[R]{MiniNumpy}
    \fancyfoot[C]{\thepage}
}

% =================================================
% 8. INFORMATIONS DU DOCUMENT
% =================================================

\title{
    \textbf{\color{blue}Le Mans Université}
    \color{black}
    \\ Licence Informatique
    \\ \textit{3ème année}
    \\ \textbf{Mini Numpy}
}

\author{
    \href{https://lounol72.github.io/Lounol72/}{Louis Subtil}
}

\date{\today}

% =================================================
% 9. DEBUT DU DOCUMENT
% =================================================

\begin{document}

% -------------------------------------------------
% PAGE DE GARDE
% -------------------------------------------------

\begin{figure}
    \includegraphics[width=0.3\textwidth]{logolemansU.png}
    \hfill
    \includegraphics[width=0.3\textwidth]{logo_IC2.png}
\end{figure}

\maketitle

% -------------------------------------------------
% TABLE DES MATIERES
% -------------------------------------------------

\newpage
\tableofcontents
\newpage

% =================================================
% CHAPITRE 1 : INTRODUCTION
% =================================================

\chapter{Introduction}

\section{Contexte}

MiniNumpy est une bibliothèque de calcul numérique
écrite de zéro en C++20, sans dépendance extérieure.
Elle reprend, à petite échelle, ce que fait NumPy
en Python : manipuler des vecteurs et des matrices
de réels et leur appliquer des opérations.

Le projet a un but final précis : réimplémenter en
C++, avec MiniNumpy à la place de NumPy, le réseau
de neurones du livre \emph{Neural Networks from
Scratch in Python}~\cite{nnfs}, puis l'entraîner sur
le jeu de données Fashion-MNIST (des images de
vêtements à classer en dix catégories).

C'est ce but qui fixe le contenu de ce carnet :
chaque notion présentée ici sert, à un moment ou
à un autre, à faire fonctionner ce réseau.

\section{Rôle de ce carnet}

Écrire une bibliothèque numérique oblige à
comprendre exactement ce que l'on calcule : une
formule mal comprise donne un code qui compile mais
renvoie de faux résultats. Ce carnet fait donc le
lien entre la théorie et le code.

Chaque notion y est présentée de la même façon :
\begin{enumerate}
    \item une \textbf{définition} précise ;
    \item un \textbf{exemple calculé à la main}, sur
          de petites valeurs, qui sert aussi de cas
          de test ;
    \item son \textbf{utilité dans le projet} :
          à quoi elle sert dans le réseau de
          neurones ;
    \item son \textbf{implémentation} : la classe
          ou la fonction de MiniNumpy qui la
          réalise.
\end{enumerate}

\section{Plan}

Le carnet suit l'ordre dans lequel la bibliothèque
est construite :
\begin{enumerate}
    \item \textbf{Vecteurs.} Les opérations de base,
          le produit scalaire et la norme. Un
          neurone calcule un produit scalaire entre
          ses entrées et ses poids.
    \item \textbf{Matrices.} Stockage, accès aux
          coefficients, transposée et produit
          matriciel. Les poids d'une couche entière
          forment une matrice, et un seul produit
          matriciel calcule la sortie de tous les
          neurones de la couche, pour tout un lot
          d'images.
    \item \textbf{Broadcasting.} La règle qui permet
          d'additionner un vecteur à chaque ligne
          d'une matrice, par exemple pour ajouter
          les biais d'une couche.
    \item \textbf{Fonctions d'activation.} ReLU et
          ses voisines, qui rendent le réseau non
          linéaire.
    \item \textbf{Softmax.} Elle transforme les
          sorties de la dernière couche en
          probabilités pour les dix classes.
    \item \textbf{Entropie croisée.} La fonction de
          perte, qui mesure l'écart entre la
          prédiction du réseau et la bonne réponse.
    \item \textbf{Rétropropagation.} Dérivées
          partielles, gradient et dérivation en
          chaîne : comment chaque poids influe sur
          l'erreur.
    \item \textbf{Descente de gradient.} Elle
          corrige les poids pas à pas pour faire
          baisser l'erreur ; c'est l'entraînement
          proprement dit.
\end{enumerate}

\section{Conventions de notation}

\begin{itemize}
    \item Les scalaires sont notés en minuscules
          italiques : $a$, $\lambda$, $x_i$.
    \item Les vecteurs sont notés en minuscules
          grasses : $\mathbf{x} \in \R^n$. Ce sont
          des vecteurs colonnes.
    \item Les matrices sont notées en majuscules :
          $W \in \R^{m \times n}$ a $m$ lignes et
          $n$ colonnes, et $w_{ij}$ est le
          coefficient de la ligne $i$ et de la
          colonne $j$.
    \item La transposée de $W$ est notée $W^\top$.
    \item En mathématiques, les indices commencent
          à $1$ : $\mathbf{x} = (x_1, \dots, x_n)$.
          En C++, ils commencent à $0$ : la
          composante $x_i$ correspond donc à
          \texttt{v[i-1]} dans le code.
\end{itemize}


% =================================================
% CHAPITRE 2 : VECTEURS
% =================================================

\chapter{Les vecteurs}

\section{Définition}

\begin{definition}[Vecteur]
Soit $n \in \N^*$. Un \emph{vecteur} de dimension $n$
est un $n$-uplet de réels, c'est-à-dire un élément
de $\R^n$ :
\[
    \vect{u} = \begin{pmatrix} u_1 \\ u_2 \\ \vdots \\ u_n \end{pmatrix},
    \qquad u_i \in \R \text{ pour tout } i \in \{1, \dots, n\}.
\]
Les réels $u_i$ sont les \emph{composantes} de
$\vect{u}$, et $n$ est sa \emph{dimension}
(son nombre de composantes).
\end{definition}

\begin{remark}
L'exposant de $\R^n$ est la dimension $n$, fixée
pour un vecteur donné. L'indice $i$, lui, parcourt
les composantes de $1$ à $n$.
\end{remark}

\begin{definition}[Vecteur nul]
Le vecteur nul de $\R^n$ est
$\vect{0} = (0, \dots, 0)^\top$ : toutes ses
composantes valent $0$.
\end{definition}

Dans toute la suite, $\vect{u}$ et $\vect{v}$ sont
deux vecteurs de $\R^n$ et $k \in \R$ est un réel.

\section{Opérations}

\subsection{Addition et soustraction}

\begin{definition}
Pour deux vecteurs \textbf{de même dimension} $n$ :
\[
    \vect{u} + \vect{v} =
    \begin{pmatrix} u_1 + v_1 \\ \vdots \\ u_n + v_n \end{pmatrix},
    \qquad
    \vect{u} - \vect{v} =
    \begin{pmatrix} u_1 - v_1 \\ \vdots \\ u_n - v_n \end{pmatrix}.
\]
\end{definition}

\begin{importantbox}
La somme d'un vecteur de $\R^n$ et d'un vecteur de
$\R^m$ avec $n \neq m$ n'est pas définie.
\end{importantbox}

\subsection{Multiplication par un réel}

\begin{definition}
\[
    k \vect{u} =
    \begin{pmatrix} k u_1 \\ \vdots \\ k u_n \end{pmatrix}.
\]
La division par un réel $k \neq 0$ est la
multiplication par $\frac{1}{k}$ :
$\dfrac{\vect{u}}{k} = \dfrac{1}{k}\,\vect{u}$.
\end{definition}

\subsection{Produit scalaire}

\begin{definition}[Forme analytique]
Pour deux vecteurs de même dimension :
\[
    \vect{u} \cdot \vect{v} = \sum_{i=1}^{n} u_i v_i
    = u_1 v_1 + u_2 v_2 + \dots + u_n v_n.
\]
Le résultat est un \textbf{réel}, pas un vecteur.
\end{definition}

\begin{proposition}[Forme géométrique]
Si $\vect{u}$ et $\vect{v}$ sont non nuls, et si
$\theta \in [0, \pi]$ est l'angle entre eux :
\[
    \vect{u} \cdot \vect{v}
    = \norm{\vect{u}}\,\norm{\vect{v}} \cos\theta.
\]
\end{proposition}

\begin{remark}
Les deux formes donnent le \emph{même} nombre. En
pratique, on calcule le produit scalaire par la forme
analytique, puis on s'en sert pour trouver l'angle :
\[
    \cos\theta =
    \frac{\vect{u} \cdot \vect{v}}{\norm{\vect{u}}\,\norm{\vect{v}}}
    \quad\text{donc}\quad
    \theta = \arccos\!\left(
    \frac{\vect{u} \cdot \vect{v}}{\norm{\vect{u}}\,\norm{\vect{v}}}
    \right).
\]
Si $\vect{u} \cdot \vect{v} = 0$, les vecteurs
sont \emph{orthogonaux} ($\theta = \frac{\pi}{2}$).
\end{remark}

\subsection{Norme}

\begin{definition}[Norme euclidienne]
\[
    \norm{\vect{u}} = \sqrt{\vect{u} \cdot \vect{u}}
    = \sqrt{\sum_{i=1}^{n} u_i^2}.
\]
C'est la longueur du vecteur. Elle est toujours
positive, et nulle seulement pour $\vect{u} = \vect{0}$.
\end{definition}

\subsection{Normalisation}

\begin{definition}
Pour $\vect{u} \neq \vect{0}$, le vecteur normalisé
(ou unitaire) de $\vect{u}$ est
\[
    \hat{\vect{u}} = \frac{1}{\norm{\vect{u}}}\,\vect{u}.
\]
Il a la même direction et le même sens que
$\vect{u}$, et sa norme vaut $1$.
\end{definition}

\begin{importantbox}
Le vecteur nul ne peut pas être normalisé : on
diviserait par $\norm{\vect{0}} = 0$.
\end{importantbox}

\section{Exemples calculés à la main}

Tous les exemples utilisent les deux vecteurs de $\R^3$
\[
    \vect{u} = \begin{pmatrix} 1 \\ 2 \\ 2 \end{pmatrix},
    \qquad
    \vect{v} = \begin{pmatrix} 2 \\ -2 \\ 1 \end{pmatrix}.
\]

\begin{example}[Addition et soustraction]
\[
    \vect{u} + \vect{v} =
    \begin{pmatrix} 1 + 2 \\ 2 + (-2) \\ 2 + 1 \end{pmatrix}
    = \begin{pmatrix} 3 \\ 0 \\ 3 \end{pmatrix},
    \qquad
    \vect{u} - \vect{v} =
    \begin{pmatrix} 1 - 2 \\ 2 - (-2) \\ 2 - 1 \end{pmatrix}
    = \begin{pmatrix} -1 \\ 4 \\ 1 \end{pmatrix}.
\]
\end{example}

\begin{example}[Multiplication et division par un réel]
Avec $k = 2$ :
\[
    2\vect{u} =
    \begin{pmatrix} 2 \times 1 \\ 2 \times 2 \\ 2 \times 2 \end{pmatrix}
    = \begin{pmatrix} 2 \\ 4 \\ 4 \end{pmatrix},
    \qquad
    \frac{\vect{u}}{2} =
    \begin{pmatrix} 0{,}5 \\ 1 \\ 1 \end{pmatrix}.
\]
\end{example}

\begin{example}[Produit scalaire analytique]
\[
    \vect{u} \cdot \vect{v}
    = 1 \times 2 + 2 \times (-2) + 2 \times 1
    = 2 - 4 + 2 = 0.
\]
Les vecteurs $\vect{u}$ et $\vect{v}$ sont donc
orthogonaux.
\end{example}

\begin{example}[Norme]
\[
    \norm{\vect{u}} = \sqrt{1^2 + 2^2 + 2^2}
    = \sqrt{9} = 3,
    \qquad
    \norm{\vect{v}} = \sqrt{2^2 + (-2)^2 + 1^2}
    = \sqrt{9} = 3.
\]
\end{example}

\begin{example}[Produit scalaire géométrique et angle]
Pour $\vect{u}$ et $\vect{v}$ :
\[
    \cos\theta = \frac{0}{3 \times 3} = 0
    \quad\text{donc}\quad
    \theta = \frac{\pi}{2},
\]
ce qui confirme l'orthogonalité trouvée plus haut.

Pour un cas où l'angle n'est pas droit, prenons
$\vect{a} = (1, 0)^\top$ et $\vect{b} = (1, 1)^\top$
dans $\R^2$ : $\vect{a} \cdot \vect{b} = 1$,
$\norm{\vect{a}} = 1$ et $\norm{\vect{b}} = \sqrt{2}$, donc
\[
    \cos\theta = \frac{1}{1 \times \sqrt{2}} = \frac{\sqrt{2}}{2}
    \quad\text{donc}\quad
    \theta = \frac{\pi}{4}.
\]
On vérifie la forme géométrique :
$\norm{\vect{a}}\,\norm{\vect{b}}\cos\theta
= 1 \times \sqrt{2} \times \frac{\sqrt{2}}{2} = 1
= \vect{a} \cdot \vect{b}$.
\end{example}

\begin{example}[Normalisation]
\[
    \hat{\vect{u}} = \frac{1}{3}\begin{pmatrix} 1 \\ 2 \\ 2 \end{pmatrix}
    = \begin{pmatrix} 1/3 \\ 2/3 \\ 2/3 \end{pmatrix},
    \qquad
    \norm{\hat{\vect{u}}}
    = \sqrt{\tfrac{1}{9} + \tfrac{4}{9} + \tfrac{4}{9}} = 1.
\]
\end{example}

\begin{example}[Cas d'erreur]
\begin{itemize}
    \item $(1, 2)^\top + (1, 2, 3)^\top$ n'est pas défini :
          les dimensions sont $2$ et $3$.
    \item $\vect{0} \in \R^3$ ne se normalise pas,
          car $\norm{\vect{0}} = 0$.
\end{itemize}
\end{example}

\section{Utilité dans le projet}

Un neurone reçoit des entrées $\vect{x} \in \R^n$,
possède un poids par entrée, rangés dans
$\vect{w} \in \R^n$, et un biais $b \in \R$. Sa sortie,
avant la fonction d'activation, est
\[
    z = \vect{w} \cdot \vect{x} + b.
\]
Le produit scalaire est donc le calcul de base de
tout le réseau. La condition « même dimension »
correspond à une exigence concrète : il faut autant de
poids que d'entrées.


\section{Implémentation}

\begin{codebox}
Voici un exemple d'implémentation en C++20.
\end{codebox}

\lstinputlisting[
language=C++,
caption=Code d'implémentation de la fonction membre "normalisation",
firstline=12,
lastline=23,
numbers=left,
breaklines=true
]{../../src/Vector.cpp}

\newpage

\chapter{Les matrices}

\section{Définition}

\begin{definition}[Matrice]
Soient $m \in \N^*$ et $n \in \N^*$. Une
\emph{matrice} à $m$ lignes et $n$ colonnes est un
tableau rectangulaire de $m \times n$ réels : il y a
$m$ lignes, et chaque ligne contient $n$ éléments.
\[
    A =
    \begin{pmatrix}
        a_{11} & a_{12} & \cdots & a_{1n} \\
        a_{21} & a_{22} & \cdots & a_{2n} \\
        \vdots & \vdots & \ddots & \vdots \\
        a_{m1} & a_{m2} & \cdots & a_{mn}
    \end{pmatrix},
    \qquad a_{ij} \in \R.
\]
Le réel $a_{ij}$ est le \emph{coefficient} de la
ligne $i$ et de la colonne $j$. On note aussi
$A = (a_{ij})_{1 \le i \le m,\, 1 \le j \le n}$.
L'ensemble des matrices à $m$ lignes et $n$ colonnes
est noté $\R^{m \times n}$, et le couple $(m, n)$
est la \emph{taille} de $A$.
\end{definition}

\begin{remark}
On donne toujours le nombre de lignes \emph{avant}
le nombre de colonnes, aussi bien pour la taille
$m \times n$ que pour l'indice $a_{ij}$. Une matrice
$2 \times 3$ et une matrice $3 \times 2$ ne sont pas
de la même taille.
\end{remark}

\begin{definition}[Lignes et colonnes]
Pour $A \in \R^{m \times n}$ :
\begin{itemize}
    \item la \emph{ligne} $i$ de $A$ est
          $(a_{i1}, a_{i2}, \dots, a_{in})$ : elle a
          $n$ éléments ;
    \item la \emph{colonne} $j$ de $A$ est
          $(a_{1j}, a_{2j}, \dots, a_{mj})^\top$ :
          elle a $m$ éléments.
\end{itemize}
\end{definition}

\begin{definition}[Égalité]
Deux matrices $A$ et $B$ sont égales si elles ont
\textbf{la même taille} et si $a_{ij} = b_{ij}$ pour
tous $i$ et $j$.
\end{definition}

\section{Matrices particulières}

\begin{definition}[Matrice carrée]
Une matrice est \emph{carrée} si elle a autant de
lignes que de colonnes ($m = n$). Ses coefficients
$a_{11}, a_{22}, \dots, a_{nn}$ forment sa
\emph{diagonale}.
\end{definition}

\begin{definition}[Matrice nulle]
La matrice nulle de $\R^{m \times n}$, notée $0$
(ou $0_{m,n}$), a tous ses coefficients égaux à $0$.
\end{definition}

\begin{definition}[Matrice identité]
La matrice identité $I_n \in \R^{n \times n}$ est la
matrice carrée qui a des $1$ sur la diagonale et des
$0$ partout ailleurs :
\[
    I_n =
    \begin{pmatrix}
        1 & 0 & \cdots & 0 \\
        0 & 1 & \cdots & 0 \\
        \vdots & \vdots & \ddots & \vdots \\
        0 & 0 & \cdots & 1
    \end{pmatrix},
    \qquad
    (I_n)_{ij} =
    \begin{cases}
        1 & \text{si } i = j, \\
        0 & \text{sinon.}
    \end{cases}
\]
\end{definition}

\begin{remark}
Un vecteur $\vect{x} \in \R^n$ peut être vu comme une
matrice à $n$ lignes et une seule colonne, c'est-à-dire
un élément de $\R^{n \times 1}$. C'est cohérent avec la
convention de ce carnet : les vecteurs sont des
vecteurs colonnes.
\end{remark}

Dans toute la suite, $A$ et $B$ sont deux matrices
de $\R^{m \times n}$ et $k \in \R$ est un réel.

\section{Opérations}

\subsection{Addition et soustraction}

\begin{definition}
Pour deux matrices \textbf{de même taille}, on
additionne (ou soustrait) coefficient par
coefficient :
\[
    (A + B)_{ij} = a_{ij} + b_{ij},
    \qquad
    (A - B)_{ij} = a_{ij} - b_{ij}.
\]
Le résultat est encore une matrice de
$\R^{m \times n}$.
\end{definition}

\begin{importantbox}
La somme de deux matrices de tailles différentes
n'est pas définie, même si elles ont le même nombre
total de coefficients (par exemple une $2 \times 3$
et une $3 \times 2$).
\end{importantbox}

\subsection{Multiplication par un réel}

\begin{definition}
\[
    (kA)_{ij} = k\, a_{ij}.
\]
Chaque coefficient est multiplié par $k$. Comme pour
les vecteurs, la division par un réel $k \neq 0$ est
la multiplication par $\frac{1}{k}$.
\end{definition}

\subsection{Transposée}

\begin{definition}[Transposée]
La \emph{transposée} de $A \in \R^{m \times n}$ est
la matrice $A^\top \in \R^{n \times m}$ définie par
\[
    (A^\top)_{ij} = a_{ji}.
\]
Les lignes de $A$ deviennent les colonnes de
$A^\top$, et inversement.
\end{definition}

\begin{importantbox}
La transposée échange les dimensions : si $A$ a
$m$ lignes et $n$ colonnes, $A^\top$ a $n$ lignes et
$m$ colonnes. Elle ne garde la même taille que pour
une matrice carrée.
\end{importantbox}

\begin{proposition}
Pour $A, B \in \R^{m \times n}$ et $k \in \R$ :
\begin{enumerate}
    \item $(A^\top)^\top = A$ ;
    \item $(A + B)^\top = A^\top + B^\top$ ;
    \item $(kA)^\top = k\, A^\top$.
\end{enumerate}
\end{proposition}

\begin{definition}[Matrice symétrique]
Une matrice carrée $A$ est \emph{symétrique} si
$A^\top = A$, c'est-à-dire si $a_{ij} = a_{ji}$ pour
tous $i$ et $j$. Par exemple, $I_n$ est symétrique.
\end{definition}

\subsection{Produit de Hadamard}

\begin{definition}[Produit terme à terme]
Le \emph{produit de Hadamard} de deux matrices
\textbf{de même taille} est la matrice
$A \odot B \in \R^{m \times n}$ définie par
\[
    (A \odot B)_{ij} = a_{ij}\, b_{ij}.
\]
\end{definition}

\begin{importantbox}
Ce n'est \emph{pas} le produit matriciel défini
ci-dessous. Attention à NumPy : l'opérateur
\texttt{*} y calcule $A \odot B$, et c'est
\texttt{@} (ou \texttt{np.dot}) qui calcule $AB$.
\end{importantbox}

\subsection{Produit matriciel}

\begin{definition}[Produit matriciel]
Soient $A \in \R^{m \times n}$ et
$B \in \R^{n \times p}$. Le \emph{produit}
$AB \in \R^{m \times p}$ est défini par
\[
    (AB)_{ij} = \sum_{k=1}^{n} a_{ik}\, b_{kj}
    = a_{i1} b_{1j} + a_{i2} b_{2j} + \dots + a_{in} b_{nj}.
\]
Le coefficient $(AB)_{ij}$ est le produit scalaire
de la ligne $i$ de $A$ et de la colonne $j$ de $B$.
\end{definition}

\begin{importantbox}
Le produit $AB$ n'existe que si le nombre de
\textbf{colonnes} de $A$ est égal au nombre de
\textbf{lignes} de $B$ :
\[
    \underbrace{(m \times n)}_{A}
    \cdot
    \underbrace{(n \times p)}_{B}
    = \underbrace{(m \times p)}_{AB}.
\]
Les deux dimensions « intérieures » doivent être
égales, et les deux dimensions « extérieures » donnent
la taille du résultat.
\end{importantbox}

\begin{proposition}
Pour des matrices dont les tailles rendent les
produits possibles, et pour $k \in \R$ :
\begin{enumerate}
    \item associativité : $(AB)C = A(BC)$ ;
    \item distributivité : $A(B + C) = AB + AC$ et
          $(A + B)C = AC + BC$ ;
    \item $(kA)B = A(kB) = k(AB)$ ;
    \item élément neutre : pour
          $A \in \R^{m \times n}$,
          $I_m A = A I_n = A$ ;
    \item transposée d'un produit :
          $(AB)^\top = B^\top A^\top$.
\end{enumerate}
\end{proposition}

\begin{importantbox}
Le produit matriciel n'est \textbf{pas commutatif} :
en général $AB \neq BA$. Même quand les deux
produits existent, ils peuvent ne pas avoir la même
taille, et même quand $A$ et $B$ sont carrées, ils
sont en général différents.
\end{importantbox}

\subsection{Produit d'une matrice par un vecteur}

\begin{definition}
Soient $A \in \R^{m \times n}$ et $\vect{x} \in \R^n$.
En voyant $\vect{x}$ comme une matrice de
$\R^{n \times 1}$, le produit $A\vect{x}$ est un
vecteur de $\R^m$ :
\[
    (A\vect{x})_i = \sum_{k=1}^{n} a_{ik}\, x_k,
    \qquad i \in \{1, \dots, m\}.
\]
La composante $i$ de $A\vect{x}$ est le produit
scalaire de la ligne $i$ de $A$ avec $\vect{x}$.
\end{definition}

\begin{remark}
Le produit scalaire est lui-même un produit
matriciel : pour $\vect{u}, \vect{v} \in \R^n$,
$\vect{u}^\top \in \R^{1 \times n}$ et
$\vect{v} \in \R^{n \times 1}$, donc
$\vect{u}^\top \vect{v} \in \R^{1 \times 1}$ est la
matrice qui contient le réel
$\vect{u} \cdot \vect{v}$.
\end{remark}

\subsection{Trace}

\begin{definition}[Trace]
La \emph{trace} d'une matrice carrée
$A \in \R^{n \times n}$ est la somme de ses
coefficients diagonaux :
\[
    \operatorname{tr}(A) = \sum_{i=1}^{n} a_{ii}.
\]
Par exemple, $\operatorname{tr}(I_n) = n$.
\end{definition}

% =================================================
% BIBLIOGRAPHIE
% =================================================

\printbibliography

\end{document}
